\subsection{挠积的定义与基本性质}

对于每个A-模E，我们用 \(p_{E}: L(E) \to E\) 表示E的典范自由分解（X，第 50 页）。

定义1.——设M为右A-模，N为左A-模。M与N的挠积是分次k-模

\[
\operatorname{Tor} ^ {\mathbf {A}} (\mathbf {M}, \mathbf {N}) = \mathrm{H} (\mathrm{L} (\mathbf {M}) \otimes_ {\mathbf {A}} \mathrm{L} (\mathbf {N}))  .\tag{4}
\]

Tor \(^{A}\) (M, N) 的齐次分量记为

\[
\operatorname{Tor} _ {n} ^ {\mathbf {A}} (\mathbf {M}, \mathbf {N}) = \mathrm{H} _ {n} (\mathrm{L} (\mathbf {M}) \otimes_ {\mathbf {A}} \mathrm{L} (\mathbf {N})).\tag{5}
\]

由于L(M)和L(N)在右侧为零，我们有

\[
\operatorname{Tor} _ {n} ^ {\mathbf {A}} (\mathbf {M}, \mathbf {N}) = 0 \quad \text { for } \quad n <   0.\tag{6}
\]

注记1. --- 我们将在下文（X，第107页，命题6）看到Tor \(^{A}\) (M, N) 模的有限性性质。例如，若 A 是交换诺特环，且 M 与 N 是有限生成的 A-模，则每个 A-模 Tor \(_{n}^{A}\) (M, N) 是有限生成的。

设 \(f: \mathbf{M} \to \mathbf{M}'\) 是右A-模的同态，且 \(g: \mathbf{N} \to \mathbf{N}'\) 为左A-模的一个同态；我们令 \(\operatorname{Tor}^{\mathbf{A}}(f, g) = \operatorname{H}\left(\operatorname{L}(f) \otimes_{\mathbf{A}} \operatorname{L}(g)\right)\) ；它是分次 \(k\) -模的一个同态

\[
\operatorname{Tor} ^ {\mathbf {A}} (f, g): \operatorname{Tor} ^ {\mathbf {A}} (\mathbf {M}, \mathbf {N}) \rightarrow \operatorname{Tor} ^ {\mathbf {A}} (\mathbf {M} ^ {\prime}, \mathbf {N} ^ {\prime})
\]

其齐次分量记为

\[
\operatorname{Tor} _ {n} ^ {\mathbf {A}} (f, g): \operatorname{Tor} _ {n} ^ {\mathbf {A}} (\mathbf {M}, \mathbf {N}) \rightarrow \operatorname{Tor} _ {n} ^ {\mathbf {A}} (\mathbf {M} ^ {\prime}, \mathbf {N} ^ {\prime}).
\]

根据X，第62页，命题1，典范同态

\[
\gamma_ {0, 0}: \mathrm{H} _ {0} (\mathrm{L} (\mathrm{M})) \otimes_ {\mathrm{A}} \mathrm{H} _ {0} (\mathrm{L} (\mathrm{N})) \rightarrow \mathrm{H} _ {0} (\mathrm{L} (\mathrm{M}) \otimes_ {\mathrm{A}} \mathrm{L} (\mathrm{N}))
\]

是双射的；利用同构 \(\mathbf{M} \to \mathbf{H}_{0}(\mathbf{L}(\mathbf{M}))\) 与 \(\mathbf{N} \to \mathbf{H}_{0}(\mathbf{L}(\mathbf{N}))\)，我们推导出一个同构，称为典范同构

\[
\gamma_ {\mathbf {M}, \mathbf {N}}: \mathbf {M} \otimes_ {\mathbf {A}} \mathbf {N} \rightarrow \operatorname{Tor} _ {0} ^ {\mathbf {A}} (\mathbf {M}, \mathbf {N}).\tag{7}
\]

我们将始终通过这个同构把 \(\operatorname{Tor}_{0}^{\mathbf{A}}(\mathbf{M},\mathbf{N})\) 与 \(\mathbf{M}\otimes_{\mathbf{A}}\mathbf{N}\) 等同起来。于是 \(k\) -线性映射 \(\operatorname{Tor}_{0}^{\mathbf{A}}(f,g)\) 与 \(f\otimes g\) 等同。

注记2. --- 复形的态射 \(p_{\mathbf{M}} \otimes p_{\mathbf{N}}: \mathbf{L}(\mathbf{M}) \otimes_{\mathbf{A}} \mathbf{L}(\mathbf{N}) \to \mathbf{M} \otimes_{\mathbf{A}} \mathbf{N}\) 在0次同调上诱导同构

\[
\gamma_ {\mathbf {M}, \mathbf {N}} ^ {- 1}: \operatorname{Tor} _ {0} ^ {\mathbf {A}} (\mathbf {M}, \mathbf {N}) \rightarrow \mathbf {M} \otimes_ {\mathbf {A}} \mathbf {N}
\]

\(\gamma_{M,N}\) 的逆。

我们有 \(L(1_{\mathbf{M}}) = 1_{L(\mathbf{M})}\) , \(L(1_{\mathbf{N}}) = i_{L(\mathbf{N})}\) 因此通过取同调：

\[
\operatorname{Tor} ^ {\mathbf {A}} \left(1 _ {\mathbf {M}}, 1 _ {\mathbf {N}}\right) = 1 _ {\operatorname{Tor} ^ {\mathbf {A}} (\mathbf {M}, \mathbf {N})}.\tag{8}
\]

若 \(f': \mathbf{M}' \to \mathbf{M}''\) (相应地， \(g': \mathbf{N}' \to \mathbf{N}''\) )是右（相应地，左）A-模的同态，我们有 \(\mathbf{L}(g' \circ g) = \mathbf{L}(g') \circ \mathbf{L}(g)\) 与 \(\mathbf{L}(f' \circ f) = \mathbf{L}(f') \circ \mathbf{L}(f)\) ，因此

\[
\operatorname{Tor} ^ {\mathbf {A}} \left(f ^ {\prime} \circ f, g ^ {\prime} \circ g\right) = \operatorname{Tor} ^ {\mathbf {A}} \left(f ^ {\prime}, g ^ {\prime}\right) \circ \operatorname{Tor} ^ {\mathbf {A}} (f, g).\tag{9}
\]

考虑k-复形的态射

\[
\mathrm{L} (\mathrm{M}) \otimes_ {\mathrm{A}} \mathrm{N} \xleftarrow {1 \otimes p _ {\mathrm{N}}} \mathrm{L} (\mathrm{M}) \otimes_ {\mathrm{A}} \mathrm{L} (\mathrm{N}) \xrightarrow {p _ {\mathrm{M}} \otimes 1} \mathrm{M} \otimes_ {\mathrm{A}} \mathrm{L} (\mathrm{N})
\]

以及它们在同调中诱导的k-同态：

\[
\mathrm{H} (\mathrm{L} (\mathrm{M}) \otimes_ {\mathrm{A}} \mathrm{N}) \xleftarrow {\psi_ {\mathrm{M}} (\mathrm{N})} \operatorname{Tor} ^ {\mathrm{A}} (\mathrm{M}, \mathrm{N}) \xrightarrow {\overline {{\psi}} _ {\mathrm{N}} (\mathrm{M})} \mathrm{H} (\mathrm{M} \otimes_ {\mathrm{A}} \mathrm{L} (\mathrm{N}));
\]

根据X，第67页，命题4， \(1 \otimes p_{\mathbf{N}}\) 与 \(p_{\mathbf{M}} \otimes 1\) 是拟同构。因此：

命题5. --- k-同态

\[
\begin{array}{l} \psi_ {\mathrm{M}} (\mathrm{N}): \operatorname{Tor} ^ {\mathrm{A}} (\mathrm{M}, \mathrm{N}) \to \mathrm{H} (\mathrm{L} (\mathrm{M}) \otimes_ {\mathrm{A}} \mathrm{N}) \\ \overline {{\psi}} _ {\mathrm{N}} (\mathrm{M}): \operatorname{Tor} ^ {\mathrm{A}} (\mathrm{M}, \mathrm{N}) \to \mathrm{H} (\mathrm{M} \otimes_ {\mathrm{A}} \mathrm{L} (\mathrm{N})) \end{array}
\]

都是双射的。

推论。—— 若 M 或 N 是平坦的，则 \(\operatorname{Tor}_{i}^{\mathbf{A}}(\mathbf{M},\mathbf{N}) = 0\) 对 \(i \geqslant 0\) 成立。

假设 N（相应地 M）是平坦的；则 \(p_{\mathbf{M}} \otimes 1 : \mathbf{L}(\mathbf{M}) \otimes_{\mathbf{A}} \mathbf{N} \to \mathbf{M} \otimes_{\mathbf{A}} \mathbf{N}\) (相应地， \(1 \otimes p_{\mathbf{N}} : \mathbf{M} \otimes_{\mathbf{A}} \mathbf{L}(\mathbf{N}) \to \mathbf{M} \otimes_{\mathbf{A}} \mathbf{N}\) ）是一个拟同构（X，第 67 页，命题 4），因此 \(\mathrm{H}_{i}(\mathrm{L}(\mathrm{M}) \otimes_{\mathbf{A}} \mathrm{N})\) (相应地， \(\mathrm{H}_{i}(\mathrm{M} \otimes_{\mathbf{A}} \mathrm{L}(\mathrm{N}))\) ）对 i \textgreater{} 0 为零。

注记 3。—— 若 \(g: \mathbf{N} \to \mathbf{N}'\) 是左 A-模的同态，则

\[
\left(1 _ {\mathrm{L} (\mathbf {M})} \otimes g\right) \circ \left(1 _ {\mathrm{L} (\mathbf {M})} \otimes 1 _ {\mathrm{N}}\right) = \left(1 _ {\mathrm{L} (\mathbf {M})} \otimes 1 _ {\mathrm{N}}\right) \circ \left(1 _ {\mathrm{L} (\mathbf {M})} \otimes \mathrm{L} (g)\right),
\]

因此图表

\[
\begin{array}{c c c} \operatorname{Tor} ^ {\mathbf {A}} (\mathbf {M}, \mathbf {N}) & \xrightarrow {\psi_ {\mathbf {M}} (\mathbf {N})} & \mathrm{H} (\mathrm{L} (\mathbf {M}) \otimes_ {\mathbf {A}} \mathbf {N}) \\ \operatorname{Tor} ^ {\mathbf {A}} (1, g) \Big \downarrow & & \Big \downarrow \mathrm{H} (1 \otimes g) \\ \operatorname{Tor} ^ {\mathbf {A}} (\mathbf {M}, \mathbf {N} ^ {\prime}) & \xrightarrow {\psi_ {\mathbf {M}} (\mathbf {N} ^ {\prime})} & \mathrm{H} (\mathrm{L} (\mathbf {M}) \otimes_ {\mathbf {A}} \mathbf {N} ^ {\prime}) \end{array}
\]

是交换的。

同样地，若 \(f: \mathbf{M} \to \mathbf{M}'\) 是右 A-模的同态，我们有交换图：

\[
\begin{array}{c c} \operatorname{Tor} ^ {\mathbf {A}} (\mathrm{M}, \mathrm{N}) & \xrightarrow {\overline {{\psi}} _ {\mathrm{N}} (\mathrm{M})} \mathrm{H} (\mathrm{M} \otimes_ {\mathrm{A}} \mathrm{L} (\mathrm{N})) \\ \operatorname{Tor} ^ {\mathbf {A}} (f, 1) \Big \downarrow & \Big \downarrow \mathrm{H} (f \otimes 1) \\ \operatorname{Tor} ^ {\mathbf {A}} (\mathrm{M} ^ {\prime}, \mathrm{N}) & \xrightarrow {\overline {{\psi}} _ {\mathrm{N}} (\mathrm{M} ^ {\prime})} \mathrm{H} (\mathrm{M} ^ {\prime} \otimes_ {\mathrm{A}} \mathrm{L} (\mathrm{N})). \end{array}
\]

命题6. --- 映射 (f, g) ↦ Tor \(^{A}\) (f, g) :

\[
\operatorname{Hom} _ {\mathrm{A}} \left(\mathrm{M}, \mathrm{M} ^ {\prime}\right) \times \operatorname{Hom} _ {\mathrm{A}} \left(\mathrm{N}, \mathrm{N} ^ {\prime}\right)\rightarrow \operatorname{Hom} _ {k} \left(\operatorname{Tor} ^ {\mathrm{A}} (\mathrm{M}, \mathrm{N}), \operatorname{Tor} ^ {\mathrm{A}} \left(\mathrm{M} ^ {\prime}, \mathrm{N} ^ {\prime}\right)\right)
\]

是 k-双线性的。

设 \(f \in \operatorname{Hom}_{\mathbf{A}}(\mathbf{M}, \mathbf{M}')\) , \(g_1, g_2 \in \operatorname{Hom}_{\mathbf{A}}(\mathbf{N}, \mathbf{N}')\) , \(\lambda_1, \lambda_2 \in k\) 。则态射

\[
\lambda_ {1} (\mathrm{L} (f) \otimes g _ {1}) + \lambda_ {2} (\mathrm{L} (f) \otimes g _ {2}) \quad \text { and } \quad \mathrm{L} (f) \otimes (\lambda_ {1} g _ {1} + \lambda_ {2} g _ {2})
\]

这两个从 L(M) ⊗\_A N 到 L(M) ⊗\_A N′ 的态射相一致；由命题5和注记3，我们因此有

\[
\operatorname{Tor} ^ {\mathbf {A}} \left(f, \lambda_ {1} g _ {1} + \lambda_ {2} g _ {2}\right) = \lambda_ {1} \operatorname{Tor} ^ {\mathbf {A}} \left(f, g _ {1}\right) + \lambda_ {2} \operatorname{Tor} ^ {\mathbf {A}} \left(f, g _ {2}\right).\tag{10}
\]

我们对映射 \(f\mapsto\mathrm{Tor}^{\mathbf{A}}(f,g)\) 作类似推理。

推论. --- 设 \(\lambda\in k\) 。若 \(\lambda\) 湮灭 M 或 N，则它湮灭 Tor \(^{A}\) (M, N)。

事实上， \(\lambda .1_{\mathrm{Tor}(\mathbf{M},\mathbf{N})} = \mathrm{Tor}(\lambda .1_{\mathbf{M}},1_{\mathbf{N}}) = \mathrm{Tor}(1_{\mathbf{M}},\lambda .1_{\mathbf{N}})\) .

命题7. --- 设I和J为两个集合， \((\mathbf{M}_{\alpha})_{\alpha \in I}\) 一族右 A-模， \((\mathbf{N}_{\beta})_{\beta \in J}\) 一个左A-模族。同态

\[
\bigoplus_ {\alpha \in I, \beta \in J} \operatorname{Tor} ^ {A} \left(M _ {\alpha}, N _ {\beta}\right)\rightarrow \operatorname{Tor} ^ {A} \left(\bigoplus_ {\alpha \in I} M _ {\alpha}, \bigoplus_ {\beta \in J} N _ {\beta}\right)
\]

由典范同态 \(M_{\alpha} \rightarrow \oplus M_{\alpha}\) 与 \(N_{\beta} \rightarrow \oplus N_{\beta}\) 导出的，是双射的。只需证明对每个右模M（相应地，每个左模N），典范同态

\[
\bigoplus_ {\beta \in J} \operatorname{Tor} ^ {A} (M, N _ {\beta}) \rightarrow \operatorname{Tor} ^ {A} (M, \bigoplus_ {\beta \in J} N _ {\beta})
\]

\((\text{resp. } \bigoplus_{\alpha \in I} \text{Tor}^{\mathbf{A}}(\mathbf{M}_{\alpha}, \mathbf{N}) \to \text{Tor}^{\mathbf{A}} (\bigoplus_{\alpha \in I} \mathbf{M}_{\alpha}, \mathbf{N}))\) 是双射的。这由前述内容、X，第28页，命题1以及以下典范同构得出：

\[
\begin{array}{l} \underset {\beta} {\oplus} (\mathrm{L(M)} \otimes_ {\mathrm{A}} \mathrm {N}_ {\beta}) \to \mathrm{L(M)} \otimes_ {\mathrm{A}} (\underset {\beta} {\oplus} \mathrm {N}_ {\beta}), \\ \underset {\alpha} {\oplus} (\mathrm {M}_ {\alpha} \otimes_ {\mathrm{A}} \mathrm{L(N)}) \to (\underset {\alpha} {\oplus} \mathrm {M}_ {\alpha}) \otimes_ {\mathrm{A}} \mathrm{L(N)}. \end{array}
\]

类似的论证给出：

命题8.——设I（相应地，J）是一个右有向预序集， \(((M_{\alpha}), (u_{\alpha' \alpha}))\) (相应地， \(((N_{\beta}), (v_{\beta' \beta}))\) ) 是关于 I（相应地，J）的右（相应地，左）A-模的归纳系。分次 k-模的同态

\[
\lim _ {(\alpha , \beta) \in I \times J} \operatorname{Tor} ^ {A} \left(M _ {\alpha}, N _ {\beta}\right)\rightarrow \operatorname{Tor} ^ {A} \left(\lim _ {\alpha \in I} M _ {\alpha}, \lim _ {\beta \in J} N _ {\beta}\right),
\]

由典范的 A-同态 \(\mathbf{M}_{\alpha} \to \varinjlim \mathbf{M}_{\alpha}\) 与 \(\mathbf{N}_{\beta} \to \varinjlim \mathbf{N}_{\beta}\) 导出的，是双射。

特别地，取 \(J = I\) 并注意到 \((\alpha, \alpha), \alpha \in I\) 构成 \(I \times I\) 的共尾子集，我们得到：

推论. --- 设 I 是一个右有向预序集， \((\mathbf{M}_{i}, u_{ji})\) (相应地， \((\mathbf{N}_{i}, v_{ji})\) ) 是关于 I 的右（相应地，左）A-模的归纳系。分次 k-模的同态

\[
\varinjlim_ {i \in I} \operatorname{Tor} ^ {A} (M _ {i}, N _ {i}) \rightarrow \operatorname{Tor} ^ {A} (\varinjlim_ {i \in I} M _ {i}, \varinjlim_ {i \in I} N _ {i}),
\]

由典范的 A-同态 \(\mathbf{M}_i\to \varinjlim \mathbf{M}_i\) 与 \(\mathbf{N}_j\to \varinjlim \mathbf{N}_j\) 导出的，是双射。

设 M 为右 A-模，N 为左 A-模，A° 为 A 的反环，M° 为 M 所对应的左 A°-模，N° 为 N 所对应的右 A°-模。我们有 L(M°) = L(M)° 且 L(N°) = L(N)°，从而得到一个交换同构（X，第 63 页，命题 2）

\[
\sigma (\mathrm{L} (\mathrm{M}), \mathrm{L} (\mathrm{N})): \mathrm{L} (\mathrm{M}) \otimes_ {\mathrm{A}} \mathrm{L} (\mathrm{N}) \rightarrow \mathrm{L} (\mathrm{N} ^ {\circ}) \otimes_ {\mathrm{A} ^ {0}} \mathrm{L} (\mathrm{M} ^ {\circ}).
\]

过渡到同调， \(\sigma(\mathrm{L}(\mathbf{M}),\mathrm{L}(\mathbf{N}))\) 诱导出一个次数为 0 的分次同构 \(\sigma_{\mathbf{M},\mathbf{N}}:\operatorname{Tor}^{\mathbf{A}}(\mathbf{M},\mathbf{N})\to\operatorname{Tor}^{\mathbf{A}^{\circ}}(\mathbf{N}^{\circ},\mathbf{M}^{\circ})\) 称为挠积的交换同构。

注意 \(\sigma_{\mathbf{N}^{\circ},\mathbf{M}^{\circ}}\circ\sigma_{\mathbf{M},\mathbf{N}}=\mathrm{Id}_{\mathrm{Tor}\left(\mathbf{M},\mathbf{N}\right)}\) 并且 \(\sigma_{\mathbf{M},\mathbf{N}}\) 在零次项上诱导出张量积的交换同态。另一方面，若 \(f:\mathbf{M}\to\mathbf{M}'\) 与 \(g:\mathbf{N}\to\mathbf{N}'\) 是A-模的同态，我们有

\[
\operatorname{Tor} ^ {\mathbf {A} ^ {\circ}} (g, f) \circ \sigma_ {\mathbf {M}, \mathbf {N}} = \sigma_ {\mathbf {M} ^ {\prime}, \mathbf {N} ^ {\prime}} \circ \operatorname{Tor} ^ {\mathbf {A}} (f, g).
\]

